\section{Combinatorial conventions and notations}\label{sec:combinatorics}

Fix an indeterminate~$q$.
We def\/ine the~$q$-numbers,~$q$-factorials, and~$q$-binomials~by
\begin{gather*}
(n)=(n)_q=1+q+\cdots+q^{n-1},
\\
(n)!=(n)\cdot(n-1)\cdots(2)\cdot(1),
\qquad
{n\choose k}_q=\frac{(n)!}{(k)!\cdot(n-k)!}.
\end{gather*}
In def\/ining our generalized Feigin homomorphism $\overline\Psi_\mathbf{i}$ we will make particular use of the
bar-invariant versions of the quantum numbers, for convenience we introduce $v=\sqrt{q}$:
\begin{gather*}
[n]=[n]_q=v^{-n+1}+v^{-n+3}+\cdots+v^{n-3}+v^{n-1},
\\
[n]!=[n]\cdot[n-1]\cdots[2]\cdot[1],
\qquad
{n\brack k}_q=\frac{[n]!}{[k]!\cdot[n-k]!}.
\end{gather*}
The~$q$-binomials satisfy several useful analogues of classical binomial identities.
We collect these and sketch their proofs in the following.
\begin{Lemma}\label{lem:binomial_identities}\qquad
\begin{enumerate}\itemsep=0pt
\item[$1.$] The~$q$-binomial coefficients satisfy the following identities:
\begin{enumerate}[\upshape (a)]\itemsep=0pt
\item Pascal identities: ${n\choose k}_q={n-1\choose k-1}_q+q^k{n-1\choose k}_q=q^{n-k}{n-1\choose k-1}_q+{n-1\choose k}_q$;
\item row identity: for $n>0$, $\sum\limits_{k=0}^n (-1)^k q^{-nk+\frac{1}{2} k(k+1)}{n\choose k}_q=0$;
\item subspace identity: ${m+n\choose k}_q=\sum\limits_{r+s=k}q^{r(n-s)}{m\choose r}_q{n\choose s}_q$.
\end{enumerate}
\item[$2.$] The bar-invariant~$q$-binomial coefficients satisfy the following identities:
\begin{enumerate}[\upshape (a)]\itemsep=0pt
\item Pascal identities: ${n\brack k}_q=v^{n-k}{n-1\brack k-1}_q+v^{-k}{n-1\brack k}_q=v^{k-n}{n-1\brack k-1}_q+v^k{n-1\brack k}_q$;
\item row identity: for $n>0$, $\sum\limits_{k=0}^n (-1)^k v^{k-nk} {n\brack k}_q=0$;
\item subspace identity: ${m+n\brack k}_q=\sum\limits_{r+s=k}v^{r(n-s)-s(m-r)}{m\brack r}_q{n\brack s}_q$.
\end{enumerate}
\end{enumerate}
\end{Lemma}

\begin{proof}
We prove the identities in (1).
The Pascal identities in (a) follow from the def\/inition and the equalities
$(n)_q=(k)_q+q^k(n-k)_q=q^{n-k}(k)_q+(n-k)_q$.
For (b) we simply apply the second Pascal identity to get a~telescopic summation:
\begin{gather*}
\sum\limits_{k=0}^n (-1)^k q^{-nk+\frac{1}{2} k(k+1)}{n\choose k}_q
\\
\qquad
=\sum\limits_{k=1}^n (-1)^k
q^{-n(k-1)+\frac{1}{2}(k-1)k}{n-1\choose k-1}_q+\sum\limits_{k=0}^{n-1} (-1)^k q^{-nk+\frac{1}{2} k(k+1)}{n-1\choose
k}_q=0.
\end{gather*}
To see (c) we consider a~f\/inite f\/ield $\mathbb{F}$ with~$q$ elements.
Write $\operatorname{Gr}_k(\mathbb{F}^n)$ for the set of~$k$-dimensional subspaces of $\mathbb{F}^n$.
The number of points in $\operatorname{Gr}_k(\mathbb{F}^n)$ is given by ${n\choose k}_q$.
Fix a~decomposition $\mathbb{F}^{m+n}=\mathbb{F}^m\oplus\mathbb{F}^n$ so that any subspace of $\mathbb{F}^{m+n}$ can be
decomposed as a~direct sum of a~subspace of $\mathbb{F}^m$ and a~subspace of $\mathbb{F}^n$.
This gives rise to a~surjective map $\operatorname{Gr}_k(\mathbb{F}^{m+n})\onto \bigsqcup\limits_{r+s=k}\operatorname{Gr}_r(\mathbb{F}^m)\times\operatorname{Gr}_s(\mathbb{F}^n)$ with f\/iber over any point
of $\operatorname{Gr}_r(\mathbb{F}^m)\times\operatorname{Gr}_s(\mathbb{F}^n)$ an af\/f\/ine space of dimension $r(n-s)$.
Now counting points completes the proof when~$q$ is a~power of a~prime.
But this is an equality of polynomials for inf\/initely many values and thus the polynomials must be equal for every~$q$.

The identities in (2) easily follow from those in (1) using the relation ${n\brack k}_q=v^{-k(n-k)}{n\choose k}_q$.
\end{proof}

We choose a~notation for symmetric groups which is most convenient for describing the structure constants of the
multiplication in the quantum shuf\/f\/le algebra.
For each positive integer~$t$ let $\Sigma_t$ denote the symmetric group on~$t$ letters, thought of as the set of all
orderings $\sigma=(\sigma_1,\ldots,\sigma_t)$ of the set $(1,\ldots,t)$.
We write $\sigma\cdot\tau=(\sigma_{\tau_1},\ldots,\sigma_{\tau_t})$ for the multiplication in $\Sigma_t$.
Let $\sigma^{-1}_k$ denote the position of~$k$ in~$\sigma$, i.e the number~$i$ so that $\sigma_i=k$.

For positive integers~$r$ and~$s$ def\/ine the \emph{shuffle subgroup} $\Sigma_{r,s}\subset\Sigma_{r+s}$ consisting of all
shuf\/f\/les of the sequences $(1,\ldots,r)$ and $(r+1,\ldots,r+s)$, that is
\begin{gather*}
\Sigma_{r,s}=\big\{\sigma\in\Sigma_{r+s}: \sigma^{-1}_1<\cdots<\sigma^{-1}_r, \sigma^{-1}_{r+1}<\cdots<\sigma^{-1}_{r+s}\big\}.
\end{gather*}
By convention we take $\Sigma_{0,0}=\Sigma_0=\{1\}$ to be the trivial group.

\section{Quantum groups and Feigin homomorphisms}\label{sec:feigin}

In this section we introduce the main objects underlying the present investigations: the quantum groups associated to
a~symmetrizable Cartan matrix and the Feigin homomorphisms mapping them to quantum polynomial algebras.

Fix an index set~$I$.
Let $A=(a_{ij})$ be an $I\times I$ Cartan matrix with symmetrizing matrix $D=\operatorname{diag}(\mathbf{d})$, where
$\mathbf{d}=(d_i: i\in I)$ is an~$I$-tuple of positive integers.
More explicitly, the entries of~$A$ have the following properties:
\begin{itemize}\itemsep=0pt
\item $a_{ii}=2$ for all~$i$;
\item $a_{ij}\le0$ for $i\ne j$;
\item $d_ia_{ij}=d_ja_{ji}$ for all $i,j\in I$.
\end{itemize}
Let $\mathcal{Q}$ denote the root lattice of a~Kac--Moody root system~$\Phi$ associated to~$A$.
Denoting by $\{\alpha_i\}_{i\in I}$ the simple roots of~$\Phi$, $\mathcal{Q}$ is the free abelian %Abelian ?
group with basis $\{\alpha_i\}$.
The pair $(A,\mathbf{d})$ determines a~symmetric bilinear form $(\cdot,\cdot):\mathcal{Q}\times\mathcal{Q}\to\mathbb{Z}$
given on generators of $\mathcal{Q}$ by $(\alpha_i,\alpha_j)=d_ia_{ij}$ for $i,j\in I$.
